% =====================================================================
%  MIC Summer-of-Code  --  Week 6
%  Deep Generative Modeling: VAEs, GANs, and a peek at Diffusion
%
%  Self-contained handbook chapter.  Compile with pdflatex (twice) or
%  upload to Overleaf and hit Recompile.
%      pdflatex MIC_VAEs_and_GANs.tex
%      pdflatex MIC_VAEs_and_GANs.tex
%
%  Built from Prof. Suyash P. Awate's CS-736 DeepGenerativeModeling deck
%  (VAE/GAN half), with attribution, for the use of mentees on this project.
% =====================================================================
\documentclass[11pt]{article}

\usepackage[a4paper,margin=1in]{geometry}
\usepackage{amsmath,amssymb,amsthm,mathtools,bm}
\usepackage{enumitem}
\usepackage{booktabs}
\usepackage{xcolor}
\usepackage{microtype}
\usepackage[most]{tcolorbox}
\usepackage{titlesec}
\usepackage[colorlinks=true,linkcolor=blue!45!black,urlcolor=blue!55!black,
            citecolor=blue!45!black]{hyperref}

\titleformat{\section}{\Large\bfseries\sffamily\color{blue!40!black}}{\thesection}{0.6em}{}
\titleformat{\subsection}{\large\bfseries\sffamily\color{blue!32!black}}{\thesubsection}{0.5em}{}

\newtcolorbox{keybox}[1][]{colback=blue!4!white,colframe=blue!40!black,
  fonttitle=\bfseries\sffamily,title=#1,boxrule=0.6pt,arc=2pt,left=4pt,right=4pt,top=3pt,bottom=3pt}
\newtcolorbox{notebox}[1][Aside]{colback=black!3!white,colframe=black!45!white,
  fonttitle=\bfseries\sffamily,title=#1,boxrule=0.5pt,arc=2pt,left=4pt,right=4pt,top=3pt,bottom=3pt}

\newcommand{\R}{\mathbb{R}}
\newcommand{\E}{\mathbb{E}}
\newcommand{\KL}{\mathrm{KL}}
\newcommand{\norm}[1]{\left\lVert #1 \right\rVert}
\newcommand{\deck}[1]{\textsf{\small[#1]}}
\DeclareMathOperator*{\argmin}{arg\,min}
\DeclareMathOperator*{\argmax}{arg\,max}

\title{\textbf{Deep Generative Models}\\[2pt]
\large VAEs, GANs, and a peek at diffusion --- concept and applications\\[6pt]
\normalsize Medical Image Computing -- Summer of Code -- Week 6}
\author{}
\date{}

\begin{document}
\maketitle
\vspace{-2.2em}

\begin{keybox}[How to read this chapter]
Last week a network \emph{classified} or \emph{segmented} an image. This week a
network \emph{generates} one. A generative model learns the distribution of
images $p(x)$ well enough to draw brand-new, realistic samples from it. We cover
the three families you will meet --- \textbf{VAEs}, \textbf{GANs}, and (in
preview) \textbf{diffusion} --- at the level of \emph{concept and applications},
not a derivation marathon. The single thread to hold onto, and the reason this
sits inside a Bayesian course at all: \textbf{a generative model is a learned
prior}, and it plugs straight into the reconstruction objective you have used
since Week 3. Keep the \deck{DeepGenerativeModeling} deck open (the generative
material is slides 33--78); do the notebook for the concepts in NumPy. Real
PyTorch/MONAI practice lives in the cited tutorials.
\end{keybox}

\tableofcontents
\bigskip

% =====================================================================
\section{Generative vs discriminative, and the core idea}
% =====================================================================

A \textbf{discriminative} model learns a boundary --- ``is this $x$ a cat or a
dog?'', i.e.\ $p(\text{label}\mid x)$. A \textbf{generative} model learns the
\emph{data itself} --- the distribution $p(x)$ of plausible images --- so that it
can both score how typical a new image is and \emph{synthesise} new ones
\deck{DeepGenerativeModeling, slide 33}.

How do you represent a distribution over something as high-dimensional as an
image? The trick behind every model this week is the same
\deck{DeepGenerativeModeling, slide 34}:

\begin{keybox}[The generative recipe]
Start from a simple distribution that is trivial to sample --- a
\textbf{standard normal} $Z\sim\mathcal{N}(0,I)$ in a low-dimensional
\emph{latent space}. Learn a neural network $f_\theta$ that maps a latent sample
to an image, $x = f_\theta(z)$, so that the induced distribution of generated
images matches the true data distribution $p(x)$. To make a new image: sample
$z$, push it through the decoder. The families differ only in \emph{how} they
train $f_\theta$ to make the match.
\end{keybox}

% =====================================================================
\section{Autoencoders: the warm-up}
% =====================================================================

Recall the encoder--decoder from Week 5. An \textbf{autoencoder} (AE) is the
unsupervised version: an encoder compresses an image $X$ to a small latent code
$h$, and a decoder reconstructs $X'\approx X$ from it
\deck{DeepGenerativeModeling, slides 27--29}. Trained to minimise reconstruction
error, the code $h$ becomes an \emph{efficient representation} of the data ---
useful for denoising, anomaly detection, and as features. Regularising the AE
(a bottleneck, weight penalties, dropout, noise) stops it from learning the
trivial copy and forces it to capture structure.

But a plain AE is \emph{not} generative: its latent space has holes, so a random
$h$ usually decodes to garbage. The fix is to give the latent space a
\emph{distribution} --- which is exactly what a VAE does.

\begin{notebox}[A linear autoencoder is PCA]
If the encoder and decoder are linear and the loss is squared error, the optimal
autoencoder spans the same subspace as \textbf{PCA} --- the latent code is the
projection onto the top principal components, and reconstruction error falls as
you keep more. So the autoencoder idea is not new; it is PCA's nonlinear,
learnable descendant. (You will see this concretely in the notebook.)
\end{notebox}

% =====================================================================
\section{Variational Autoencoders (VAEs)}
% =====================================================================

A \textbf{VAE} makes the autoencoder generative by turning the latent code into a
\emph{random variable} \deck{DeepGenerativeModeling, slides 49--51}. Instead of
mapping $X$ to a single point $h$, the encoder maps $X$ to a whole
\textbf{distribution} $q_\phi(Z\mid X)$ over latents (it outputs a mean and a
variance). The decoder $p_\theta(X\mid Z)$ maps a latent sample back to an image.
``Variational'' refers to this modelling of a \emph{family} of latent codes per
input, which captures uncertainty and --- crucially --- fills in the latent space
so that sampling works.

\subsection{The objective: the ELBO}

We would like to maximise the data likelihood $\log p(x)$, but the integral over
latents is intractable. The standard move is to maximise a tractable lower bound,
the \textbf{Evidence Lower BOund} \deck{DeepGenerativeModeling, slides 52--55}:
\[
\log p(x) \;\ge\; \underbrace{\E_{q_\phi(z\mid x)}\!\big[\log p_\theta(x\mid z)\big]}_{\text{reconstruction}}
\;-\; \underbrace{\KL\!\big(q_\phi(z\mid x)\,\|\,p(z)\big)}_{\text{regulariser}} .
\]
Read it as a tug-of-war:
\begin{itemize}[nosep,leftmargin=*]
  \item the \textbf{reconstruction} term wants the decoded image to look like the
  input (with a Gaussian decoder this becomes the familiar squared error);
  \item the \textbf{KL} term pulls the encoder's $q_\phi(z\mid x)$ toward the
  prior $p(z)=\mathcal{N}(0,I)$ --- it must not wander off, so that latents we
  sample at generation time decode to something sensible. But it must not collapse
  \emph{onto} the prior either, or the code carries no information about $x$.
\end{itemize}
Design choices make it concrete: $p(z)=\mathcal{N}(0,I)$ (easy to sample at
generation), encoder outputs $(\mu_\phi(x),\sigma_\phi(x))$, decoder Gaussian with
fixed variance \deck{DeepGenerativeModeling, slide 54}.

\begin{notebox}[Where the ELBO comes from]
The bound is not a guess. For \emph{any} encoder $q_\phi(z\mid x)$,
\[
\log p(x)=\underbrace{\E_{q_\phi}\!\Big[\log\tfrac{p_\theta(x,z)}{q_\phi(z\mid x)}\Big]}_{\text{ELBO}}
   \;+\; \KL\!\big(q_\phi(z\mid x)\,\|\,p(z\mid x)\big).
\]
The second term is a KL divergence, so it is $\ge 0$; hence $\log p(x)\ge$ ELBO,
with equality exactly when the encoder matches the true posterior $p(z\mid x)$.
Maximising the ELBO therefore does double duty: it pushes up the data likelihood
\emph{and} pulls the approximate posterior toward the true one.
\end{notebox}

\begin{notebox}[Worked example: KL between two Gaussians]
The KL term has a closed form, which is why VAEs are practical. For a diagonal
Gaussian posterior $\mathcal{N}(\mu,\sigma^2)$ against the standard normal,
per latent dimension,
\[
\KL\!\big(\mathcal{N}(\mu,\sigma^2)\,\|\,\mathcal{N}(0,1)\big)
   = \tfrac12\big(\sigma^2 + \mu^2 - 1 - \log\sigma^2\big).
\]
Sum over dimensions and you have the entire regulariser --- no sampling needed.
\end{notebox}

\subsection{The reparameterisation trick}

There is a snag: training needs gradients to flow back through the \emph{random}
sampling $z\sim q_\phi(z\mid x)$, and you cannot backpropagate through a random
draw. The \textbf{reparameterisation trick} moves the randomness out of the way
\deck{DeepGenerativeModeling, slide 56}:
\[
z = \mu_\phi(x) + \sigma_\phi(x)\odot\varepsilon, \qquad \varepsilon\sim\mathcal{N}(0,I).
\]
Now $z$ is a \emph{deterministic, differentiable} function of the encoder outputs
$(\mu_\phi,\sigma_\phi)$ plus external noise $\varepsilon$, so gradients pass
straight through to $\phi$. It is a small trick with outsized importance --- it is
what makes VAEs trainable by ordinary backprop.

\subsection{What VAEs are good at}

VAEs give a \emph{smooth, structured} latent space: interpolating between two
codes produces sensible in-between images, and sampling $z\sim\mathcal{N}(0,I)$
generates new ones. In medical imaging they shine at \textbf{anomaly detection}
(train only on healthy scans; a lesion reconstructs poorly, so high
reconstruction error flags it), \textbf{denoising} (a VAE returns a whole
\emph{family} of clean candidates, not one), \textbf{data augmentation}, and
representation learning \deck{DeepGenerativeModeling, slides 57--58}. Their
weakness: samples tend to be a little \emph{blurry}, because the Gaussian
reconstruction loss averages over plausible outputs.

% =====================================================================
\section{Generative Adversarial Networks (GANs)}
% =====================================================================

A \textbf{GAN} learns to generate by staging a contest between two networks
\deck{DeepGenerativeModeling, slides 38--43}:
\begin{itemize}[nosep,leftmargin=*]
  \item the \textbf{generator} $G$ maps a latent $z\sim\mathcal{N}(0,I)$ to a fake
  image;
  \item the \textbf{discriminator} $D$ tries to tell real images from $G$'s fakes,
  outputting $D(\text{img})\to1$ for ``real'', $\to0$ for ``fake''.
\end{itemize}
They are trained against each other in a \textbf{min--max game}:
\[
\min_G \max_D\;\; \E_{x\sim p_{\text{data}}}\!\big[\log D(x)\big]
   + \E_{z}\!\big[\log\big(1 - D(G(z))\big)\big].
\]
$D$ improves at spotting fakes; $G$ improves at fooling $D$; at the ideal
equilibrium $G$'s samples are indistinguishable from real data. A short analysis
shows that for a \emph{fixed} generator the optimal discriminator is
$D^\star(x)=\dfrac{p_{\text{data}}(x)}{p_{\text{data}}(x)+p_g(x)}$, and plugging it
back reveals that $G$ is minimising the \textbf{Jensen--Shannon divergence}
between the generated and real distributions \deck{DeepGenerativeModeling, slides
40--41}. No explicit likelihood, no reconstruction loss --- just ``make samples a
critic cannot distinguish''.

\paragraph{Why GANs are hard.} The adversarial game is a saddle-point problem,
not a simple minimisation, and it can misbehave \deck{DeepGenerativeModeling,
slides 42--44}:
\begin{itemize}[nosep,leftmargin=*]
  \item \textbf{mode collapse} --- $G$ finds a few images that fool $D$ and emits
  only those, losing diversity;
  \item \textbf{vanishing gradients} --- if $D$ gets too good too fast, $G$ gets no
  useful signal;
  \item \textbf{non-convergence} --- the two can oscillate forever.
\end{itemize}
The standard stabilisers: \textbf{WGAN} / \textbf{WGAN-GP} (replace JS with the
Wasserstein distance plus a gradient penalty), \textbf{spectral normalisation}
on $D$, two-time-scale learning rates (TTUR), and EMA of the generator weights.

\paragraph{Conditional GANs.} Feed both networks a condition $y$ ($G(z,y)$,
$D(x,y)$) and you can \emph{control} the output: \textbf{pix2pix} (paired
image-to-image translation) and \textbf{CycleGAN} (unpaired) are the workhorses
for medical \textbf{modality translation} --- e.g.\ synthesising a CT from an MRI
\deck{DeepGenerativeModeling, slide 45}. GANs give the sharpest samples of the
three families, at the cost of the trickiest training.

% =====================================================================
\section{Diffusion models (a preview)}
% =====================================================================

\textbf{Diffusion models} are today's state of the art for image synthesis, and
the deck frames them elegantly as a \emph{hierarchical VAE}
\deck{DeepGenerativeModeling, slides 63--76}. Two processes:
\begin{itemize}[nosep,leftmargin=*]
  \item a fixed \textbf{forward process} that, over many steps, scales the signal
  down and adds Gaussian noise, turning any image $x_0$ into pure noise $x_T\approx
  \mathcal{N}(0,I)$ --- \emph{no} network involved;
  \item a learned \textbf{reverse process} that undoes one noising step at a time,
  a Gaussian $p_\theta(x_{t-1}\mid x_t)$ whose mean a network predicts.
\end{itemize}
The training objective is again an ELBO; the punchline is that the network's job
reduces to \textbf{denoising} --- predict the clean $x_0$ (equivalently, the noise)
from a noisy $x_t$, then take the analytically-exact small step back
\deck{DeepGenerativeModeling, slide 76}. Generation runs the reverse chain from
noise. Diffusion overtook GANs because it is \emph{stable to train} and produces
\emph{high-quality, diverse} samples; the cost is slow, many-step sampling. It is
also the most direct bridge back to this course: a diffusion model is literally a
trained denoiser, the very Week-2 object, stacked across noise levels.

\begin{notebox}[The one closed form that makes diffusion practical]
Because each forward step just scales the signal and adds Gaussian noise, the
composition of $t$ steps is \emph{also} Gaussian in closed form:
$q(x_t\mid x_0)=\mathcal{N}\!\big(\sqrt{\bar\alpha_t}\,x_0,\,(1-\bar\alpha_t)I\big)$
with $\bar\alpha_t=\prod_{s\le t}\alpha_s$. So during training you can jump
straight to any noise level $t$ in one shot (no $t$-step simulation), and as
$\bar\alpha_T\to0$ the image becomes pure noise --- which is why generation can
start from $\mathcal{N}(0,I)$ and walk the reverse chain back to an image.
\end{notebox}

\subsection{Controllable generation}

An unconditional model draws \emph{some} plausible image; clinically you usually
want a \emph{specific} one. \textbf{Conditioning} feeds a label, a segmentation,
or a low-quality scan to the generator --- conditional GANs ($G(z,y)$, $D(x,y)$),
class- or image-conditioned diffusion --- and \textbf{classifier-free guidance}
lets a diffusion model trade diversity for fidelity to the condition by mixing its
conditional and unconditional predictions. This is what turns ``generate a brain
MRI'' into ``generate the high-resolution MRI consistent with \emph{this}
low-resolution one'': super-resolution, inpainting, and modality translation are
all just \emph{conditional} sampling.

% =====================================================================
\section{The unifying thread: a generative model is a learned prior}
% =====================================================================

Step back and the whole course rhymes. In Week 2 the prior $p(x)$ was a Gibbs/MRF
energy you designed by hand; in Week 3 you plugged it into reconstruction,
$\hat x=\argmin_x \tfrac{1}{2\sigma^2}\norm{Gx-y}^2 + \lambda R(x)$. A VAE, GAN,
or diffusion model is simply a \textbf{learned} $p(x)$, so it can replace the
hand-built regulariser $R$:
\begin{itemize}[nosep,leftmargin=*]
  \item a VAE's negative log-likelihood as $R(x)$;
  \item a GAN/normalising-flow density as $R(x)$;
  \item a diffusion model's \emph{score} (gradient of $\log p$) plugged straight
  into the gradient step.
\end{itemize}
This is the \textbf{plug-and-play} / \textbf{score-based reconstruction} family,
currently the hottest area in accelerated-MRI and low-dose-CT reconstruction
\deck{DeepGenerativeModeling, slides 77--78}. Same objective, learned prior.

\begin{notebox}[Which family when?]
\begin{tabular}{@{}lll@{}}
\toprule
Model & Strength & Typical medical use \\
\midrule
Autoencoder & simple, fast & denoising, features \\
VAE & smooth latent, principled & anomaly detection, augmentation \\
GAN & sharpest samples & modality translation, synthesis \\
Diffusion & best quality, stable & synthesis, super-res, recon priors \\
\bottomrule
\end{tabular}
\end{notebox}

% =====================================================================
\section{Applications --- and the trust question}
% =====================================================================

Generative models address three structural problems in medical imaging
\deck{DeepGenerativeModeling, slide 45}: \emph{scarce labelled data}
(augmentation, synthetic training sets for rare conditions), \emph{privacy}
(share synthetic instead of real scans), and \emph{anomaly detection} (flag what
the model finds improbable). Concrete uses: cross-modality synthesis (MRI$\to$CT),
super-resolution of low-dose scans, inpainting, and learned priors for
accelerated reconstruction.

\begin{notebox}[Be honest about synthetic data]
Generative power cuts both ways. GANs can \textbf{hallucinate} a realistic-looking
lesion that was never there, or \emph{erase} a real one --- a clinical hazard, not
a cosmetic flaw. Diffusion models can produce chest X-rays expert radiologists
cannot distinguish from real, which is wonderful for augmentation and alarming for
provenance. So generated images are evaluated with care: distribution-level
metrics (\textbf{FID}, MS-SSIM) for realism, \emph{downstream task} performance
for usefulness, and explicit checks for \emph{privacy leakage}. The rule of
thumb: never trust a synthetic image to add information that was not in the data.
\end{notebox}

\subsection{How do you evaluate a generated image?}

You cannot eyeball a thousand synthetic scans, and ``looks real'' is not a
number. Three kinds of metric are used. \textbf{Distribution realism:}
\textbf{FID} (Frechet Inception Distance) compares the feature statistics of the
real and generated sets (lower is better), while \textbf{MS-SSIM} probes sample
diversity (to catch mode collapse). \textbf{Downstream utility:} train a model on
synthetic data and test it on \emph{real} data --- if performance holds, the
samples carry genuine signal. \textbf{Privacy:} verify the generator has not
memorised and reproduced individual training patients (a real medico-legal risk).
The honest summary: a generative model is worth only as much as the task it
\emph{provably} improves, measured on real held-out data.

\begin{notebox}[Common pitfalls with generative models]
\begin{itemize}[nosep,leftmargin=*]
  \item \textbf{Mode collapse (GANs):} the generator emits a few ``safe'' samples;
  measure diversity, not just sharpness.
  \item \textbf{Posterior collapse (VAEs):} the KL term wins and the latent is
  ignored; anneal the KL weight or strengthen the decoder.
  \item \textbf{Blurry VAE samples:} the Gaussian reconstruction loss averages
  plausible outputs; a sharper likelihood or a VQ/adversarial term helps.
  \item \textbf{Evaluation theatre:} a cherry-picked sample grid proves nothing ---
  report FID / MS-SSIM and a downstream-task number on real data.
\end{itemize}
\end{notebox}

\begin{keybox}[Do the notebook --- and finish Project 2]
\texttt{notebooks/01\_generative\_models\_concepts.ipynb} (NumPy + Matplotlib)
makes the ideas concrete without a GPU: a linear autoencoder $=$ PCA
(reconstruct + a latent ``spectrum''), \textbf{generative sampling and
interpolation} in a Gaussian latent space, and a small 2D \textbf{adversarial}
demo showing a generator distribution chase a target while a discriminator tries
to separate them. For real VAEs/GANs/diffusion, the MONAI Generative Models
tutorials are the on-ramp. This week also hosts \textbf{Project 2} (brain-MRI
segmentation, see \texttt{project\_segmentation.md}).
\end{keybox}

% =====================================================================
\section*{Resources (watch one, read one)}
\addcontentsline{toc}{section}{Resources}
% =====================================================================

\textbf{VAEs}
\begin{itemize}[nosep,leftmargin=*]
  \item Kingma \& Welling (2013), \href{https://arxiv.org/abs/1312.6114}{\emph{Auto-Encoding Variational Bayes}} (the paper); Lilian Weng, \href{https://lilianweng.github.io/posts/2018-08-12-vae/}{\emph{From Autoencoder to Beta-VAE}}.
  \item Arxiv Insights, \href{https://www.youtube.com/watch?v=9zKuYvjFFS8}{\emph{Variational Autoencoders}} (video); \href{https://www.statlect.com/fundamentals-of-probability/Kullback-Leibler-divergence}{KL closed form}.
\end{itemize}

\textbf{GANs}
\begin{itemize}[nosep,leftmargin=*]
  \item Goodfellow et al.\ (2014), \href{https://arxiv.org/abs/1406.2661}{\emph{Generative Adversarial Nets}}; Lilian Weng, \href{https://lilianweng.github.io/posts/2017-08-20-gan/}{\emph{From GAN to WGAN}}.
  \item Gulrajani et al.\ (2017), \href{https://arxiv.org/abs/1704.00028}{\emph{Improved Training of Wasserstein GANs}} (WGAN-GP).
\end{itemize}

\textbf{Diffusion}
\begin{itemize}[nosep,leftmargin=*]
  \item Ho, Jain \& Abbeel (2020), \href{https://arxiv.org/abs/2006.11239}{\emph{Denoising Diffusion Probabilistic Models}}; Lilian Weng, \href{https://lilianweng.github.io/posts/2021-07-11-diffusion-models/}{\emph{What are Diffusion Models?}}.
\end{itemize}

\textbf{Medical generative modelling \& hands-on}
\begin{itemize}[nosep,leftmargin=*]
  \item \href{https://github.com/Project-MONAI/GenerativeModels}{MONAI Generative Models} and the \href{https://github.com/Project-MONAI/tutorials/tree/main/generation}{generation tutorials} (VAEs, VQ-GANs, latent diffusion on 2D/3D scans).
  \item \href{https://book.clinicalai.guide/chapters/07-generative-models.html}{The Clinical AI Field Guide --- Generative Models} (AE/VAE/GAN/diffusion with the trust question); \href{https://arxiv.org/pdf/2410.17664}{\emph{Deep Generative Models for 3D Medical Image Synthesis}} (2024 review).
  \item Reconstruction priors: Jalal et al.\ (2021), \href{https://arxiv.org/abs/2108.01368}{\emph{Robust CS-MRI with Deep Generative Priors}}; Chung \& Ye (2022), \href{https://www.sciencedirect.com/science/article/pii/S1361841522001268}{\emph{Score-based diffusion for accelerated MRI}}.
\end{itemize}

\vfill
\begin{center}\small\itshape
Built for the MIC Summer-of-Code from Prof.\ Suyash P.\ Awate's CS-736
DeepGenerativeModeling deck (VAE/GAN half), with attribution. For mentee use on
this project; please do not redistribute the bundled decks outside it.
\end{center}

\end{document}
